\section{Results}

The results are organized to evaluate the proposed optimization framework under varying planning objectives. The analysis begins with deployments based solely on node utilities, illustrating how alternative weight combinations reflect different goals such as demand coverage, multimodal integration, or social equity. The framework is then extended by incorporating network-level metrics, demonstrating how trade-offs between proximity and accessibility reshape the resulting station layouts. The role of slope is subsequently examined through altitude-adjusted distances, which better represent the effort of cycling in hilly areas. Finally, the framework’s applicability to system growth is illustrated through a case study of Barcelona’s l'Eixample district, where new stations are optimized while accounting for the existing network of \textit{Bicing}.

    \subsection[BSS deployments based solely on node utilities]{\ac{bss} deployments based solely on node utilities}

    To evaluate the behavior of the node‐utility function \(U_{i}\) (Equation \ref{eq:BSS_utility_no_metrics}) under different planning objectives, three representative scenarios have been selected as illustrative examples. A buffer radius \(Y = 300\,\mathrm{m}\), minimum inter‐station separation \(X = 300\,\mathrm{m}\), and total station count \(k = 250\) are held constant; only the weight vectors vary. Table~\ref{tab:scenario-defs} defines the weights for each scenario and Appendix~\ref{app_paper1:scenarios} displays the spatial distributions of the variables considered in each scenario. The first scenario (S1) intends to maximize the overall ridership by placing stations in areas of high residential density coupled with a high total count of \acp{poi}, thus capturing both trip origins and attractors. The second scenario (S2) aims to promote multimodal integration by favoring sites within easy walking distance of existing \ac{pt} stops (bus, metro, tram) and directly connected by bike lanes to locations with abundant \acp{poi}. Finally, the third scenario (S3) promotes social equity by targeting neighborhoods with lower incomes, higher unemployment, and lower educational attainment.

    \begin{table}[!htpb]
        \caption[Definition of planning scenarios (S1–S3)]{
            \textbf{Definition of planning scenarios (S1–S3).} 
            Each weight vector $W^f$ sums to 1.
        }
        \label{tab:scenario-defs}
        \footnotesize
        \centering
        \begin{tabular}{p{0.32\linewidth} p{0.68\linewidth}}
        \toprule
        \textbf{Scenario}  &  \textbf{$W_f$}\\
        \midrule
        Maximum demand coverage (S1)& \(\{\,w^{\mathrm{population}}:0.4,\;w^{\mathrm{pois\_count}}:0.35,\;w^{\mathrm{pois\_diversity}}:0.25\}\)\\
            First–Last kilometer (S2)& \(\{\,w^{\mathrm{bus}}:0.2,\;w^{\mathrm{metro}}:0.2,\;w^{\mathrm{tram}}:0.1,\;w^{\mathrm{bike\_lane}}:0.4,\;w^{\mathrm{pois\_count}}:0.1\}\)\\
            Social equity (S3)& \(\{\,w^{\mathrm{income}}:0.4,\;w^{\mathrm{unemp}}:0.3,\;w^{\mathrm{edu\_primary}}:0.3\}\)\\
        \bottomrule
        \end{tabular}
    \end{table}

    Figure \ref{fig:BSS_no_graph_metric} presents the station-placement results under the three scenarios without incorporating the network-level metrics $S_{pro}$ or $S_{acc}$. In all cases, the optimization behaves as expected, successfully prioritizing nodes aligned with the goals of each scenario. In the \textit{maximum demand coverage} scenario (S1), stations concentrate in locations with very high residential density and abundant \acp{poi}. This is reflected in S1's radar chart, which shows that on average ($\pm$ standard deviation), selected nodes exhibit a population score of $0.81 \pm 0.06$, \ac{poi} count of $0.68 \pm 0.12$, and \ac{poi} diversity of $0.64 \pm 0.16$, the highest among all scenarios. In contrast, the \textit{first–last kilometer} scenario (S2) shifts station selection toward locations near bus, metro, and tram stops, as well as dense bike lanes, resulting in higher average \ac{pt} and bike‐lane scores. Average scores increase to $0.74 \pm 0.13$ for bus, $0.40 \pm 0.21$ for metro, and $0.36 \pm 0.19$ for bike lanes, all notably higher than in S1, while also maintaining a high \ac{poi} count of $0.72 \pm 0.10$. However, this comes at the cost of lower population $0.66 \pm 0.15$ and \acp{poi} diversity $0.52 \pm 0.11$ values relative to the first scenario. 

    \begin{figure}[!htbp]
        \centering
        \includegraphics[width=1\linewidth]{0_plots/paper1/main/BSS_all_scenarios.png}
        \caption[Station-level utilities and BSS optimizations for 250 stations across scenarios S1-S3]{
            \textbf{Station-level utilities and \ac{bss} optimizations for 250 stations across scenarios S1-S3.} 
            \textbf{(A)} Maps of normalized utility scores $U_i$ for all candidate nodes under three scenarios, each defined by a different vector of factor weights. Node utilities have been min–max scaled to $[0,1]$ after weighting to facilitate direct comparison across scenarios. 
            \textbf{(B)} Station locations selected by optimizing solely on $U_i$ (no network‐based objective). Each map includes a radar‐chart inset in which the black line traces the input factor weights, the colored line shows the mean normalized factor values of the chosen nodes, and the shaded band represents $\pm1$ standard deviation of those values. The concentric grid lines on each radar chart correspond to values of 0, 0.33, 0.66, and 1. Variables correspond to the following local factors: bus stops (bus), metro lines (metro), and tram lines (tram); people with only primary education (edu); average income (income); bike lanes (lanes); \acp{poi} diversity (pois\_div); total count of \acp{poi}	 (pois\_count); population (pop); and unemployment rate (unemp). More details on the variables are provided in Table~\ref{table:local_factors_nodes}.
        }
        \label{fig:BSS_no_graph_metric}
    \end{figure}

    Finally, the \textit{social equity} scenario (S3) focuses on disadvantaged regions, which is reflected in the high values of unemployment $0.75 \pm 0.09$, primary‐education share $0.82 \pm 0.08$, and lower income $0.65 \pm 0.13$. This confirms that its weights have shifted the optimization towards areas of greater social need. Yet, when compared to S1, the socioeconomic values of S3 improve only marginally. For instance, unemployment and education already reach $0.74 \pm 0.08$ and $0.75 \pm 0.10$ in S1, despite the absence of social equity weights. This occurs because regions with many unemployed residents and low education or low income tend to coincide geographically with areas of high population, with correlations of 0.96 and 0.93 respectively. In other words, by emphasizing population in S1, many of the same disadvantaged regions that S3 targets are already being captured.

    Another important effect comes from the 300 m minimum inter-station distance constraint. In scenarios S2 and S3, where there are small regions with a lot of high score nodes, the minimum inter-station distance constraint prevents the \ac{ga} from placing multiple stations within these areas. Instead, stations are spread over larger regions where normalized node scores, while still high, are closer to citywide averages. This means that very localized, high‐score regions could be relatively under-served if they fall within the minimum separation distance. 

    Finally, the tight clustering of selected stations under the three scenarios demonstrates that relying solely on local factor scores can produce an overly compact network. If the planning goal demands a more evenly distributed and accessible layout, it is essential to incorporate network‐level considerations alongside the node‐level utility. 

    \subsection[BSS deployments with integrated node and network utilities]{\ac{bss} deployments with integrated node and network utilities}

    When network-level metrics are incorporated into the optimization  (Equation \ref{eq:full_formula}) the resulting station sets can exhibit markedly different spatial structures. This effect is illustrated in Figure~\ref{fig:alpha_screening_s1}, which shows how the S1 scenario is reshaped under varying values of~$\alpha$. Corresponding results for S2 and S3 are also provided in Appendix~\ref{app_paper1:network_scenarios}. As seen in Figure~\ref{fig:alpha_screening_s1}A, increasing~$\alpha$ towards~1 leads the algorithm to prioritize proximity between stations, whereas decreasing it towards~0 shifts the focus towards system-wide accessibility. However, setting $\alpha$ to its extreme values (0 or 1) does not guarantee a purely proximity- or accessibility-optimal solution, because the overall objective $U_{BSS}^{*}$ always retains a local-suitability component. In other words, even when $\alpha= 1$ (nominally maximum station proximity), each candidate node must still meet a minimum value of local factors (e.g. population and \acp{poi} density). If one truly wanted a pure proximity or accessibility solution, the local-suitability term could be neutralized by assigning every node the same nonzero normalized utility; in that case, $\alpha=1$ or $0$ would indeed optimize only the chosen network metric.

    \begin{figure}[!htbp]
        \centering
        \includegraphics[width=1\linewidth]{0_plots/paper1/main/alpha_screening_s1.png}
        \caption[Impact of network-metrics trade-offs on station selection for scenario S1]{
            \textbf{Impact of network-metrics trade-offs on station selection for scenario S1.}  All panels use the same S1 input weights to illustrate how modifying $\alpha$ shifts the \ac{bss} network and the node utilities $U_i$. 
            \textbf{(A)} Maps of selected stations under six conditions: a baseline optimization without network metrics, and optimizations with $\alpha$ = 0, 0.25, 0.50, 0.75, and 1.
            \textbf{(B)} displays proximity score $S_{pro}$ and accessibility score $S_{acc}$ versus $\alpha$.
            \textbf{(C)} \ac{bss} utility $U_{BSS}$ alongside box-plots of individual station utilities for each $\alpha$.
            \textbf{(D)} Change in mean local factor values of the selected stations relative to the baseline without network metrics.
        }
        \label{fig:alpha_screening_s1}
    \end{figure}

    Under scenario S1, stations selected without any network‐level adjustment are already tightly clustered, as all high‐utility nodes lie in close proximity. Consequently, using $\alpha = 1$ produces only marginal gains in stations' proximity (Figure~\ref{fig:alpha_screening_s1}B). By contrast, setting $\alpha$ = 0 substantially boosts $S_{acc}$ from 0.75 to 0.91. As a result, nodes with lower utility scores are chosen (Figure~\ref{fig:alpha_screening_s1}C). In this case, prioritizing accessibility comes at the expense of population and \ac{poi} count, however, the \acp{poi} diversity is less affected (Figure~\ref{fig:alpha_screening_s1}D). Moreover, as previously noted, some factors are spatially correlated, so optimizing for accessibility can indirectly influence variables not explicitly weighted in the objective function, most notably, unemployment and educational attainment. Conversely, when prioritizing station proximity for this scenario, the algorithm tends to select locations with better cycling infrastructure, resulting in an increase in bike lane coverage among the selected stations.

    \subsection{Altitude adjustment}

    \begin{figure}[htbp!]
        \centering
        \includegraphics[width=1\linewidth]{0_plots/paper1/main/altitude_comparison_s3.png}    \caption[Impact of altitude-adjusted distances on station selection for scenario S3.]{
        \textbf{Impact of altitude-adjusted distances on station selection for scenario S3.} 
        \textbf{(A)} Original edge distances vs. altitude-adjusted distances, colored by slope of the between origin and destination. 
        \textbf{(B)} ECDF of normalized node scores for different $\alpha$ values under the S3 scenario, with and without altitude adjustments; the baseline corresponds to the optimization without network metrics for S3. 
        \textbf{(C)} Composite utility $U^*_{BSS}$ for the resulting station sets, showing a slight trade-off when elevation is considered. 
        \textbf{(D)} Spatial distribution of normalized node scores in the study area. 
        \textbf{(E–G)} Differences in selected station locations (grid cell counts) between configurations with and without altitude adjustment under the \textbf{(E)} S3 scenario, and its network-aware versions using 
        \textbf{(F)}~$\alpha=0$ and \textbf{(G)}~$\alpha=1$. Grids cells are 500 m by 500 m. Positive values indicate more stations were selected in a given cell when altitude was considered. Background contour lines represent elevation levels.
        }
        \label{fig:altitude_s3}
    \end{figure}

    To assess how altitude-based distance adjustments influence the station selection process, Figure~\ref{fig:altitude_s3} presents results for scenario~S3 under three different optimization strategies (Appendix~\ref{app_paper1:results_altitude} shows the corresponding figures for S1 and S2). This scenario is particularly illustrative, as it features distinct clusters of high-scoring nodes situated in steep terrain. Figure~\ref{fig:altitude_s3}A shows how original edge distances $d(i,j)$ are transformed into equivalent flat distances \(d_{\mathrm{eq}}(i,j)\) by penalizing uphill segments and rewarding downhill ones. These adjustments, clamped to the interval \([-0.10, 0.10]\) to avoid extreme effects, reshape the shortest-path distance matrix and alter the resulting station configurations.

    Across all strategies, incorporating altitude tends to increase distances in steep areas, which relaxes the minimum separation constraint locally and enables more stations to be placed within high-slope, high-score clusters. This generally results in a rightward shift in the node score distribution (Figures~\ref{fig:altitude_s3}B and~\ref{fig:altitude_s3}C), as more top-performing nodes can be included. These shifts are particularly evident in steep districts such as Horta-Guinardó, Nou Barris, Sant Martí, and the Montjuïc area (Figures~\ref{fig:altitude_s3}E–G), where station counts rise in the altitude-aware solutions.

    The effects that the altitude adjustment generates on the station placement varies depending on the optimization objective. In proximity-focused configurations ($\alpha=1$), the objective favors compact clusters. Altitude adjustments slightly expand the spacing within steep regions, allowing a few additional high-score nodes to be chosen as stations without violating the minimum separation constraint. This results in marginal gains in node scores while preserving overall spatial compactness.

    In accessibility-focused configurations ($\alpha=0$), slope-adjusted distances increase the perceived remoteness of steep areas. When these regions also contain high-score nodes, (as in S1 and S3) the algorithm assigns more stations there to improve network-wide accessibility, often at the expense of flatter, low-score areas. This trade-off is visible in the spatial maps and is evident in the left tail of the ECDF, where altitude-aware solutions include fewer low-score nodes. In contrast, when high scores are located in flatter areas (as in S2), the altitude adjustment has a more modest impact, leading to less pronounced redistribution. This demonstrates that while the underlying mechanism is consistent, its effect depends strongly on the spatial correlation between slope and node utility.

    Together, these results highlight a consistent effect: adjusting distances for slope enables better exploitation of steep, high-score areas that would otherwise be constrained by proximity limits or undervalued in flat-distance metrics. The nature and extent of this effect, however, depend on whether the optimization emphasizes station's proximity, system accessibility, or local suitability alone.

    \subsection[BSS expansion]{\ac{bss} expansion}

    The proposed method can easily be applied to generate the expansion of an existing \ac{bss}. This functionality is illustrated in l'Eixample district of Barcelona under scenario~S2, which emphasizes first--last kilometer connectivity. To simulate a realistic expansion process, all 108 \textit{Bicing} stations within the district were temporarily removed and replaced with an equal number of newly selected stations, while maintaining the 410 existing stations outside the district. The optimization was carried out with $\alpha=0.5$, balancing proximity between stations and system-wide accessibility. L'Eixample district was selected due to its dense \ac{pt} infrastructure and high concentration of \acp{poi} (Appendix~\ref{app_paper1:results_bss_expansion}), making it a suitable context for evaluating the methodology.

    As shown in Figure~\ref{fig:BSS_no_graph_metric}A, l'Eixample district contains many high-scoring candidate nodes under scenario~S2. The optimized stations selected in this expansion achieve an average score of $0.76 \pm 0.08$ (range $0.63$–$1.0$), indicating that the algorithm consistently prioritizes top-scoring locations while maintaining a relatively narrow spread of values. The newly selected stations complement the spatial distribution of the existing \textit{Bicing} network while strictly respecting the minimum distance constraint, both among themselves and in relation to the fixed stations outside the district. Moreover, although the existing \textit{Bicing} stations that serve as base for the expand may include stations closer than the imposed minimum inter-station distance, the framework is designed to be capable of generating feasible solutions that comply with this constraint.  

    Several of the optimized stations are located in highly significant areas of the district. Figures~\ref{fig:results_eixample}B--E present close-up views of a selection of these locations and detailed local factor values for the chosen stations of these subareas are summarized in Table~\ref{tab:subarea_factors}. Stations in Figure~\ref{fig:results_eixample}B are situated near Plaça d'Espanya, a major multimodal hub surrounded by green spaces and retail-oriented \acp{poi} such as the Arenas shopping mall. Here, the three optimized stations combine very high bus connectivity (normalized value of $0.93$), moderate metro accessibility ($0.46$), and strong \acp{poi} density ($0.73$), with average node scores around $0.81$. Figure~\ref{fig:results_eixample}C shows Plaça de Catalunya, one of the central locations of the city, characterized by an exceptional concentration of \acp{poi} and \ac{pt} connections. The five selected stations in this area achieve the highest scores in the district, with an average node score of $0.97$, reflecting very high levels of bus and metro accessibility ($0.97$ and $0.88$, respectively) and very dense \ac{poi} coverage ($0.82$), which capture the area’s role as a focal point of commerce and tourism. Figure~\ref{fig:results_eixample}D depicts Glòries, a redeveloped area that combines strong metro and tram connectivity with a mix of cultural and commercial \acp{poi}, including the Mercat dels Encants and the Catalan National Theater. The optimized stations here record a mean node score of $0.76$, supported by the maximum tram accessibility of the BSS ($1.0$) together with balanced values for bus ($0.57$), metro ($0.28$), \acp{poi} ($0.68$), and bike lanes ($0.40$). Finally, Figure~\ref{fig:results_eixample}E focuses on Avinguda Diagonal, a major avenue lined with economic and retail-related amenities. The selected stations achieve high \ac{poi} coverage ($0.85$), good accessibility to bus ($0.66$) and metro ($0.50$), and moderate integration with the cycling network ($0.46$), yielding node scores around $0.81$. In some cases (such as in Figure~\ref{fig:results_eixample}E), selected stations may appear geographically too close to one another, even though they satisfy the minimum distance constraint. This is because distance calculations are based on the directed graph, which represents the actual navigable paths that bicycles must follow. As a result, stations that seem close in straight-line terms may still be sufficiently distant in graph-based terms, as traveling between them can require a longer detour to comply with traffic regulations and use designated cycling infrastructure. This highlights the importance of understanding how the underlying graph is constructed and how distances are defined when applying spatial constraints.

    \begin{figure}[!htbp]
        \centering
        \includegraphics[width=1\linewidth]{0_plots/paper1/main/Bicing_ampliation_good_nodes.png}  
        \caption[BSS expansion in Barcelona's l'Eixample district under scenario S2]{
            \textbf{\ac{bss} expansion in Barcelona's l'Eixample district under scenario S2.} 
            \textbf{(A)} Selected locations for the \textit{Bicing} system expansion within the district, whose boundaries are delineated by the black line. Existing \textit{Bicing} stations are shown as red stars. The optimization was performed with $\alpha$=0.5, balancing proximity and accessibility objectives. The color of each selected station reflects its normalized local score based on the local factors. Dashed rectangles indicate some illustrative areas. 
            \textbf{(B-E)} Close-up views of the highlighted areas display the spatial context around selected stations, including \ac{pt} and various categories of \acp{poi}. Each subpanel corresponds to a specific rectangle marked in panel (A).
        }    
        \label{fig:results_eixample}
    \end{figure}

    \begin{table}[!htbp]
        \caption[Normalized node scores and local factor values for expansion stations in four subareas of l’Eixample under scenario S2]{
            \textbf{Normalized node scores ($U_i$) and factor values for expansion stations in four subareas of l’Eixample under scenario S2}          
            The subareas correspond to the close-up views shown in Figure~\ref{fig:results_eixample}. Values are reported as mean $\pm$ standard deviation. "Plç." stands for Plaça, and "Av." stands for Avinguda.
        }
        \scriptsize
        \centering
        \begin{tabular}{lccccccc}
            \toprule
            Location & Stations count & $U_i$ & Bus lines & Metro lines & Tram lines & POIs total & Bike lanes \\
            \midrule
            Plç. Espanya (Fig.~\ref{fig:results_eixample}B) & 3 & $0.81 \pm 0.07$ & $0.93 \pm 0.10$ & $0.46 \pm 0.19$ & $0.00 \pm 0.00$ & $0.73 \pm 0.01$ & $0.46 \pm 0.05$ \\
            Plç. Catalunya (Fig.~\ref{fig:results_eixample}C) & 5 & $0.97 \pm 0.04$ & $0.97 \pm 0.01$ & $0.88 \pm 0.14$ & $0.00 \pm 0.00$ & $0.82 \pm 0.04$ & $0.32 \pm 0.14$ \\
            Glòries (Fig.~\ref{fig:results_eixample}D) & 2 & $0.76 \pm 0.02$ & $0.57 \pm 0.04$ & $0.28 \pm 0.00$ & $1.00 \pm 0.00$ & $0.68 \pm 0.01$ & $0.40 \pm 0.07$ \\
            Av. Diagonal (Fig.~\ref{fig:results_eixample}E) & 4 & $0.81 \pm 0.01$ & $0.66 \pm 0.00$ & $0.50 \pm 0.07$ & $0.00 \pm 0.00$ & $0.85 \pm 0.01$ & $0.46 \pm 0.06$ \\
            \bottomrule
        \end{tabular}
        \label{tab:subarea_factors}
    \end{table}


    To further evaluate the expansion results, the optimized stations were directly compared with the existing \textit{Bicing} network in l’Eixample (Figure~\ref{fig:results_eixample_comparison} and Table~\ref{tab:bicing_comparison}). Although the current \textit{Bicing} stations were not placed through an explicit optimization under scenario~S2, their performance can still be assessed with the same node score. On this basis, \textit{Bicing} stations achieve an average of $0.72 \pm 0.11$ (range $0.50$–$1.0$), while the optimized set reaches $0.76 \pm 0.08$ (range $0.63$–$1.0$). This reflects a systematic shift toward higher-scoring nodes and a narrower spread of values. At the same time, the similarity in the score distributions suggests that the real-world \textit{Bicing} implementation implicitly incorporates many of the contextual factors considered in scenario~S2. Figure~\ref{fig:results_eixample_comparison} illustrates these points: Figure~\ref{fig:results_eixample_comparison}A shows that, despite only seven stations display exact overlaps, several more optimized and \textit{Bicing} stations are located on or near the same street intersection; Figure~\ref{fig:results_eixample_comparison}B reveals that most cells differ by no more than one or two stations, confirming broadly similar spatial patterns; and Figure~\ref{fig:results_eixample_comparison}C and Table~\ref{tab:bicing_comparison} indicate that the main differences in the scores lie in higher tram accessibility (+173\%), metro accessibility (+17\%), and bike-lane integration (+8\%) for the optimized stations, while bus accessibility and \acp{poi} density remain comparable. Moreover, a closer look at the few cells with larger differences between optimized and \textit{Bicing} stations shows that the number of \acp{poi} is the most influential factor, as shown in Figure~\ref{fig:pois_and_pt_eixample}. Overall, the optimized network delivers slightly higher node scores while reproducing much of the structure of the existing system.

    \begin{figure}[!htbp]
        \centering
        \includegraphics[width=\textwidth]{0_plots/paper1/main/BSS_comparison.png}
        \caption[Comparison between optimized and actual \textit{Bicing} stations in l'Eixample]{
            \textbf{Comparison between optimized and actual \textit{Bicing} stations in l'Eixample.}
            \textbf{(A)} Overlay of optimized (yellow), actual (red), and common (blue) stations. 
            \textbf{(B)} Grid heatmap showing station count differences (optimized minus actual). 
            \textbf{(C)} Violin plots comparing normalized values of key factors across the two sets of stations.
        }
        \label{fig:results_eixample_comparison}
    \end{figure}

    \begin{table}[!htbp]
        \centering
        \footnotesize
        \caption[Comparison between optimized and actual \textit{Bicing} stations in l'Eixample]{
            \textbf{Comparison between optimized and actual \textit{Bicing} stations in l'Eixample.}
            The table reports node utilities and local factor values (mean $\pm$ standard deviation) for the expansion stations selected under scenario S2 and for the existing \textit{Bicing} stations in the district.
        }
        \label{tab:bicing_comparison}
        \begin{tabular}{lcccccc}
            \toprule
            & Bus lines & Metro lines & Tram lines & POIs total & Bike lanes & Node utility \\
            \midrule
            Actual stations & $0.73 \pm 0.14$ & $0.31 \pm 0.24$ & $0.03 \pm 0.12$ & $0.72 \pm 0.07$ & $0.46 \pm 0.17$ & $0.72 \pm 0.11$ \\
            Optimized stations & $0.76 \pm 0.12$ & $0.36 \pm 0.21$ & $0.09 \pm 0.23$ & $0.73 \pm 0.06$ & $0.50 \pm 0.16$ & $0.76 \pm 0.08$ \\
            \bottomrule
        \end{tabular}
    \end{table}